\documentclass{ximera}
\title{To infinity and beyond}
\begin{document}
\begin{abstract}
Limits as $x$ goes to infinity or negative infinity, and horizontal asymptotes. Textbook §2.6.
\end{abstract}
\maketitle

Although infinity is not a number, we can look at limits as a function gets really really big, i.e., when the function goes toward infinity.

\begin{example}
\label{example:3-22}
% work: function read off the graph below (f(1)=1, f(2)=1/4, always positive)
Let $f(x)=\answer{\frac{1}{x^{2}}}$ and notice that this function is always positive. Here is a graph of the function:
\begin{image}
\includegraphics[width=0.8\textwidth]{b48a5324-93dc-45ee-b7c2-dc0b5706bda6-041_623_883_1431_619}
\end{image}
What this means is that the higher the value of $x$, the closer to $\answer{0}$ we are going to get. Intuitively we say that the limit of $f(x)$ as $x$ goes to infinity is 0.

We write this as:
\[
\lim _{x \rightarrow \infty} f(x)=\answer{0}
\]
\end{example}

We can play this game with any function. If $f$ is a function defined on $(a, \infty)$ then
\[
\lim _{x \rightarrow \infty} f(x)=L
\]
means that the value of $f(x)$ approaches $L$ as $x$ gets larger and larger.

Similarly, you can do the same thing for negative infinity.
\[
\lim _{x \rightarrow-\infty} f(x)=L
\]

\begin{example}
\label{example:3-23}
% work: function read off the graph below (value 1 at x=0, zeros at multiples of pi, decaying oscillation)
Let $f(x)=\answer{\frac{\sin(x)}{x}}$ with the graph:
\begin{image}
\includegraphics[width=0.8\textwidth]{b48a5324-93dc-45ee-b7c2-dc0b5706bda6-042_672_947_1152_587}
\end{image}
What is
\[
\lim _{x \rightarrow-\infty} f(x)=\answer{0}
\]
\end{example}

The line $y=L$ is called a \dfn{horizontal asymptote} of the curve $f(x)$ if either
\[
\lim _{x \rightarrow \infty} f(x)=L \text { or } \lim _{x \rightarrow-\infty} f(x)=L
\]

\begin{warning}
There is a battle in mathematics on the definition of a horizontal asymptote! The majority give the definition above, but some authors additionally require that the line $f(x)$ doesn't cross the asymptote!
\end{warning}

% AUTHOR CHOICE: example chosen during conversion, change freely
\begin{exercise}
\label{exercise:3-24}
Let $f(x)=\dfrac{3x^{2}-1}{x^{2}+2}$. With a partner, find
% work: divide top and bottom by x^2: (3 - 1/x^2)/(1 + 2/x^2) -> 3 as x -> +-infinity
\[
\lim _{x \rightarrow \infty} f(x)=\answer{3} \text { and } \lim _{x \rightarrow-\infty} f(x)=\answer{3}
\]
\end{exercise}

\end{document}
